% ============ 第2章 §2.7 Hardy空间的边界值 (II)(原书 pp.141-144) ============
\section{Hardy 空间的边界值 (II)}\label{sec:2.7}

2.5 节证明了属于 $H^{p}$ 的函数 $f$ 几乎处处存在径向极限 $\lim_{r \to 1}f(re^{i\theta}) = f(e^{i\theta})$,且 $f \in L^{p}$.$H^{p}$ 在 2.1 节是在 $D$ 的内部定义的,本节证明 $H^{p}$ 可以用径向极限来定义.另外,由此证明在 $L^{p}$ 的距离下 $\lim_{r \to 1}f(re^{i\theta}) = f(e^{i\theta})$.

\begin{theorem}\label{thm:2.7.1}
设 $0 < p \leq \infty$.若 $p \neq \infty$,则
\begin{align*}
\sup_{0 \leq r < 1}M_{p}(r, f) & = \lim_{r \to 1}M_{p}(r, f) \\
& = \left\{\frac{1}{2\pi}\int_{-\pi}^{\pi}|f(e^{i\theta})|^{p}\dd\theta\right\}^{1/p} = \|f\|_{p},
\end{align*}
又
\begin{equation*}
\sup_{0 \leq r < 1}M_{\infty}(r, f) = \esssup_{-\pi \leq \theta < \pi}|f(e^{i\theta})| = \|f\|_{\infty}.
\end{equation*}
因此,$H^{p}$ 在 $1 \leq p \leq \infty$ 时是范数空间,在 $0 < p < 1$ 时是度量空间.
\end{theorem}

\begin{proof}
若 $f(z) = \sum_{j=0}^{\infty}a_{j}z^{j} \in H^{2}$,则
\begin{equation*}
M_{2}(r, f) = \left(\sum_{j = 0}^{\infty}|a_{j}|^{2}r^{2j}\right)^{1/2},
\end{equation*}
故由径向极限的存在性与 Lebesgue 积分的 Fatou 引理,
\begin{align*}
\sup_{0 \leq r < 1}M_{2}(r, f) & = \lim_{r \to 1}M_{2}(r, f) \\
& = \left(\sum_{j = 0}^{\infty}|a_{j}|^{2}\right)^{1/2} = \left\{\frac{1}{2\pi}\int_{-\pi}^{\pi}|f(e^{i\theta})|^{2}\dd\theta\right\}^{1/2}.
\end{align*}
若 $f \in H^{p}$($0 < p < \infty$),由定理~\ref{thm:2.6.4} 可分解 $f = Bg$.由于 $g$ 在 $D$ 内无零点,$[g(z)]^{p/2} \in H^{2}$,且 $|B(e^{i\theta})| = 1$ a.e. $\theta$,所以
\begin{align*}
\int_{-\pi}^{\pi}|f(re^{i\theta})|^{p}\dd\theta &\leq \int_{-\pi}^{\pi}|g(re^{i\theta})|^{p}\dd\theta \\
&\rightarrow \int_{-\pi}^{\pi}|g(e^{i\theta})|^{p}\dd\theta = \int_{-\pi}^{\pi}|f(e^{i\theta})|^{p}\dd\theta.
\end{align*}
由引理~\ref{lem:2.1.1},$\sup_{0 \leq r < 1}M_{p}(r, f) = \lim_{r \to 1}M_{p}(r, f)$,由 Fatou 引理,$p \neq \infty$ 的情形得证.当 $p = \infty$ 时,$\|f\|_{\infty} \leq \sup_{0 \leq r < 1}M_{\infty}(r, f)$ 由径向极限的存在性与最大值原理是显然的.若 $f \in H^{\infty}$ 则 $f \in H^{2}$,所以证明它可以作 Poisson 积分表示是容易的,由此导出 $\sup_{0 \leq r < 1}M_{\infty}(r, f) \leq \|f\|_{\infty}$.
\end{proof}

\begin{lemma}\label{lem:2.7.2}
设 $0 < p < \infty$ 且 $k_{n} \in L^{p}$($n = 1, 2, \ldots$).若 $\lim_{n \to \infty}k_{n}(\theta) = k(\theta)$ a.e. $\theta$,且
\begin{equation*}
\lim_{n \to \infty}\int_{-\pi}^{\pi}|k_{n}(\theta)|^{p}\dd\theta = \int_{-\pi}^{\pi}|k(\theta)|^{p}\dd\theta < \infty,
\end{equation*}
则
\begin{equation*}
\lim_{n \to \infty}\int_{-\pi}^{\pi}|k_{n}(\theta) - k(\theta)|^{p}\dd\theta = 0.
\end{equation*}
\end{lemma}

\begin{proof}
设 $E$ 为 $[-\pi, \pi]$ 的 Lebesgue 可测集,令
\begin{equation*}
\Phi_{n}(E) = \int_{E}|k_{n}(\theta)|^{p}\dd\theta, \quad \Phi(E) = \int_{E}|k(\theta)|^{p}\dd\theta.
\end{equation*}
设 $E^{c}$ 为 $E$ 在 $[-\pi, \pi]$ 中的补集,由 Fatou 引理与假设,
\begin{align*}
\Phi(E) & \leq \varliminf_{n \to \infty}\Phi_{n}(E) \leq \varlimsup_{n \to \infty}\Phi_{n}(E) \\
& = \varliminf_{n \to \infty}\Phi_{n}([-\pi, \pi]) - \varliminf_{n \to \infty}\Phi_{n}(E^{c}) \\
& \leq \Phi([-\pi, \pi]) - \Phi(E^{c}) = \Phi(E).
\end{align*}
因此,对每个 $E \subset [-\pi, \pi]$,$\lim_{n \to \infty}\Phi_{n}(E) = \Phi(E)$ 成立.对任意 $\varepsilon > 0$,取 $F \subset [-\pi, \pi]$ 满足 $\Phi(F^{c}) < \varepsilon$.取 $\delta > 0$ 使得若 $Q \subset F$ 且 $d\theta(Q) < \delta$ 则 $\Phi(Q) < \varepsilon$.这由 $\Phi$ 关于 $d\theta$ 绝对连续是可能的.由 Egorov 定理,存在 $Q \subset F$ 使 $d\theta(Q) < \delta$,且 $k_{n}$ 在 $E = F \cap Q^{c}$ 上一致收敛.因此,利用当 $a \geq 0$ 且 $b \geq 0$ 时 $(a + b)^{p} \leq 2^{p}(a^{p} + b^{p})$ 成立,
\begin{align*}
\int_{-\pi}^{\pi}|k_{n}(\theta) - k(\theta)|^{p}\dd\theta&= \int_{F^{c}}|k_{n}(\theta) - k(\theta)|^{p}\dd\theta + \int_{Q}|k_{n}(\theta) - k(\theta)|^{p}\dd\theta  \\
&\quad + \int_{E}|k_{n}(\theta) - k(\theta)|^{p}\dd\theta \\
&\leq 2^{p}\left(\Phi_{n}(F^{c}) + \Phi(F^{c}) + \Phi_{n}(Q) + \Phi(Q)\right) \\
&\quad + \int_{E}|k_{n}(\theta) - k(\theta)|^{p}\dd\theta.
\end{align*}
由 $\lim_{n \to \infty}\Phi_{n}(F^{c}) = \Phi(F^{c}) < \varepsilon$ 且 $\lim_{n \to \infty}\Phi_{n}(Q) = \Phi(Q) < \varepsilon$,以及一致收敛给出
\begin{equation*}
\lim_{n \to \infty}\int_{E}|k_{n}(\theta) - k(\theta)|^{p}\dd\theta = 0,
\end{equation*}
故对任意 $\varepsilon > 0$,
\begin{equation*}
\varlimsup_{n \to \infty}\int_{-\pi}^{\pi}|k_{n}(\theta) - k(\theta)|^{p}\dd\theta \leq 2^{p}\cdot 4\varepsilon.
\end{equation*}
\end{proof}

\begin{theorem}\label{thm:2.7.3}
当 $0 < p < \infty$ 时,若 $f \in H^{p}$,则
\begin{equation*}
\lim_{r \to 1}\int_{-\pi}^{\pi}|f(re^{i\theta}) - f(e^{i\theta})|^{p}\dd\theta = 0.
\end{equation*}
\end{theorem}

\begin{proof}
由定理~\ref{thm:2.5.3} 与定理~\ref{thm:2.7.1},$\lim_{r \to 1}f(re^{i\theta}) = f(e^{i\theta})$ a.e. $\theta$,且
\begin{equation*}
\lim_{r \to 1}\int_{-\pi}^{\pi}|f(re^{i\theta})|^{p}\dd\theta = \int_{-\pi}^{\pi}|f(e^{i\theta})|^{p}\dd\theta.
\end{equation*}
由引理~\ref{lem:2.7.2},
\begin{equation*}
\lim_{r \to 1}\int_{-\pi}^{\pi}|f(re^{i\theta}) - f(e^{i\theta})|^{p}\dd\theta = 0.
\end{equation*}
\end{proof}

\begin{corollary}\label{cor:2.7.4}
当 $0 < p \leq \infty$ 时,若 $f \in H^{p}$,则
\begin{equation*}
\lim_{r \to 1}\int_{-\pi}^{\pi}\left|\log^{+}|f(re^{i\theta})| - \log^{+}|f(e^{i\theta})|\right|\dd\theta = 0.
\end{equation*}
\end{corollary}

\begin{proof}
若 $a \geq 0, b \geq 0$ 且 $0 < p \leq 1$,则
\begin{equation*}
\left|\log^{+}a - \log^{+}b\right| \leq \frac{1}{p}|a - b|^{p},
\end{equation*}
由此与定理~\ref{thm:2.7.3} 可以证明系.上面的不等式由
\begin{equation*}
\log x \leq \frac{1}{p}(x - 1)^{p} \quad (x \geq 1)
\end{equation*}
当 $1 \leq b < a$ 且 $x = a/b$ 时可得.对 $1 < p \leq \infty$,由于若 $p < p'$ 则 $H^{p} \supset H^{p'}$,比 $0 < p \leq 1$ 的情形更显然.
\end{proof}
